% ============================================================
%  CHAPTER 7: SIMPLE HARMONIC MOTION
%  The Physics Codex: A Complete University Foundation
%  Korede Samuel Omotosho (Falcon)
%  Aurikrex Learning Systems, 2026
%
%  File: chapters/part1/ch07_simple_harmonic_motion.tex
% ============================================================

\chapter{Simple Harmonic Motion}
\label{ch:shm}

\chapterquote{Nature is not only queerer than we suppose,
but queerer than we can suppose.}{J.B.S. Haldane}

\begin{objectivebox}
By the end of this chapter, you should be able to:
\begin{enumerate}[noitemsep]
  \item Define simple harmonic motion and state its
  conditions
  \item Derive and apply the equations for displacement,
  velocity and acceleration in SHM
  \item Sketch and interpret graphs of displacement,
  velocity and acceleration against time
  \item Calculate the period and frequency of a simple
  pendulum and a mass-spring system
  \item Apply energy conservation to SHM problems
  \item Distinguish between free, damped and forced
  oscillations
  \item Explain resonance and give real examples
\end{enumerate}
\end{objectivebox}

\vspace{0.4cm}

\minitoc

\vspace{0.5cm}

% ============================================================
\section{What is Oscillatory Motion?}
% ============================================================

Look around you and you will find oscillation everywhere.
A tree branch swaying in the wind. A child on a swing.
The pendulum of a grandfather clock. The strings of a
guitar. The atoms in every solid vibrating about their
equilibrium positions. Your own heartbeat. All of these
are examples of \textbf{oscillatory motion} --- motion
that repeats itself in time about a fixed position.

The simplest and most important type of oscillatory
motion is \textbf{Simple Harmonic Motion} (SHM). It is
the foundation of wave theory, acoustics, optics,
quantum mechanics, and electrical circuits. Understanding
SHM is not optional in Physics --- it is fundamental.

% ============================================================
\section{Defining Simple Harmonic Motion}
% ============================================================

\begin{definitionbox}[Simple Harmonic Motion]
\textbf{Simple Harmonic Motion} (SHM) is a type of
oscillatory motion in which the acceleration of the
body is:
\begin{enumerate}[noitemsep]
  \item Directly proportional to the displacement
  from the equilibrium position
  \item Always directed toward the equilibrium
  position (opposite to displacement)
\end{enumerate}
\[
  a = -\omega^2 x
\]
where $x$ is displacement from equilibrium and
$\omega$ is the angular frequency. The negative sign
confirms the acceleration always opposes displacement.
\end{definitionbox}

\begin{keybox}[The Restoring Force]
For SHM to occur, there must be a \textbf{restoring force}
--- a force that always acts toward the equilibrium
position and is proportional to displacement:
\[
  F = -kx
\]
where $k$ is the restoring force constant. From
$F = ma$: $a = -kx/m$, so $\omega^2 = k/m$.
\end{keybox}

\begin{historybox}[Galileo and the Pendulum]
According to legend, in 1583 the 19-year-old Galileo
sat in Pisa Cathedral watching a chandelier swinging.
He noticed that no matter how wide or narrow the swing,
each oscillation took the same time (he timed it with
his own pulse). This property --- period independent
of amplitude --- became one of the defining features
of SHM. Galileo went on to propose using pendulums
as timekeepers, an idea Huygens turned into the first
pendulum clock in 1656.
\end{historybox}

% ============================================================
\section{Equations of Simple Harmonic Motion}
% ============================================================

The defining equation $a = -\omega^2 x$ is a
differential equation whose solution gives displacement
as a sinusoidal function of time.

\begin{lawbox}[Equations of SHM]
For a body oscillating with angular frequency $\omega$
and amplitude $A$:

\textbf{Displacement} (starting from equilibrium):
\[
  x = A\sin(\omega t) \quad \text{or} \quad
  x = A\cos(\omega t)
\]
(Use sine if $x=0$ at $t=0$; cosine if $x=A$ at $t=0$)

\textbf{Velocity:}
\[
  v = A\omega\cos(\omega t)
  = \omega\sqrt{A^2 - x^2}
\]

\textbf{Acceleration:}
\[
  a = -A\omega^2\sin(\omega t) = -\omega^2 x
\]

\textbf{Maximum values:}
\[
  v_{max} = A\omega \quad \text{(at equilibrium, } x=0\text{)}
\]
\[
  a_{max} = A\omega^2 \quad \text{(at amplitude, } x=\pm A\text{)}
\]
\end{lawbox}

\begin{diagbox}[SHM Graphs --- Displacement, Velocity, Acceleration]
\begin{center}
\begin{tikzpicture}
\begin{axis}[
  name=plot1,
  width=12cm, height=3.5cm,
  xlabel={}, ylabel={$x$},
  xmin=0, xmax=6.28, ymin=-1.3, ymax=1.3,
  xtick={0,1.571,3.142,4.712,6.283},
  xticklabels={$0$,$T/4$,$T/2$,$3T/4$,$T$},
  ytick={-1,0,1},
  yticklabels={$-A$,$0$,$A$},
  grid=major, grid style={gray!20},
  title={Displacement: $x = A\sin(\omega t)$}
]
  \addplot[codexblue, very thick, domain=0:6.28,
    samples=100]{sin(deg(x))};
\end{axis}

\begin{axis}[
  name=plot2,
  at={(plot1.below south west)},
  anchor=above north west,
  yshift=-0.3cm,
  width=12cm, height=3.5cm,
  xlabel={}, ylabel={$v$},
  xmin=0, xmax=6.28, ymin=-1.3, ymax=1.3,
  xtick={0,1.571,3.142,4.712,6.283},
  xticklabels={$0$,$T/4$,$T/2$,$3T/4$,$T$},
  ytick={-1,0,1},
  yticklabels={$-A\omega$,$0$,$A\omega$},
  grid=major, grid style={gray!20},
  title={Velocity: $v = A\omega\cos(\omega t)$}
]
  \addplot[codexred, very thick, domain=0:6.28,
    samples=100]{cos(deg(x))};
\end{axis}

\begin{axis}[
  name=plot3,
  at={(plot2.below south west)},
  anchor=above north west,
  yshift=-0.3cm,
  width=12cm, height=3.5cm,
  xlabel={Time}, ylabel={$a$},
  xmin=0, xmax=6.28, ymin=-1.3, ymax=1.3,
  xtick={0,1.571,3.142,4.712,6.283},
  xticklabels={$0$,$T/4$,$T/2$,$3T/4$,$T$},
  ytick={-1,0,1},
  yticklabels={$-A\omega^2$,$0$,$A\omega^2$},
  grid=major, grid style={gray!20},
  title={Acceleration: $a = -A\omega^2\sin(\omega t)$}
]
  \addplot[codexgreen, very thick, domain=0:6.28,
    samples=100]{-sin(deg(x))};
\end{axis}
\end{tikzpicture}
\end{center}
\small Note: acceleration is always $180°$ out of phase
with displacement.
\end{diagbox}

\begin{examplebox}[7.1]
\stepcounter{workedexample}
A particle undergoes SHM with amplitude $0.08\ \text{m}$
and period $2.0\ \text{s}$.\\
(a) Find the angular frequency.\\
(b) Find the maximum speed.\\
(c) Find the maximum acceleration.\\
(d) Find the speed and acceleration when
$x = 0.05\ \text{m}$.
\end{examplebox}

\begin{solutionbox}
\textbf{(a)} Angular frequency:
\[
  \omega = \frac{2\pi}{T} = \frac{2\pi}{2.0}
  = \pi\ \text{rad\,s}^{-1} \approx 3.14\ \text{rad\,s}^{-1}
\]

\textbf{(b)} Maximum speed:
\[
  v_{max} = A\omega = 0.08 \times \pi
  = 0.08\pi \approx 0.251\ \text{m\,s}^{-1}
\]

\textbf{(c)} Maximum acceleration:
\[
  a_{max} = A\omega^2 = 0.08 \times \pi^2
  = 0.08\pi^2 \approx 0.789\ \text{m\,s}^{-2}
\]

\textbf{(d)} At $x = 0.05\ \text{m}$:
\[
  v = \omega\sqrt{A^2-x^2}
  = \pi\sqrt{0.08^2-0.05^2}
  = \pi\sqrt{0.0064-0.0025}
\]
\[
  v = \pi\sqrt{0.0039}
  = \pi \times 0.0624
  = 0.196\ \text{m\,s}^{-1}
\]
\[
  a = -\omega^2 x = -\pi^2 \times 0.05
  = -0.493\ \text{m\,s}^{-2}
\]
The negative sign confirms acceleration is directed
back toward equilibrium (opposite to positive $x$).
\end{solutionbox}

% ============================================================
\section{The Simple Pendulum}
% ============================================================

\begin{processbox}[Derivation of Pendulum Period]
Consider a bob of mass $m$ on a string of length $L$
displaced by angle $\theta$ from vertical.

The restoring force (tangential component of gravity):
\[
  F = -mg\sin\theta
\]

For small angles ($\theta < 10°$), $\sin\theta \approx \theta$
(in radians), and the arc length $s = L\theta$:
\[
  F \approx -mg\theta = -mg\frac{s}{L}
  = -\left(\frac{mg}{L}\right)s
\]

This has the form $F = -ks$ with $k = mg/L$. Therefore
the pendulum undergoes SHM with:
\[
  \omega^2 = \frac{k}{m} = \frac{g}{L}
\]
\[
  \boxed{T = 2\pi\sqrt{\frac{L}{g}}}
\]
\end{processbox}

\begin{keybox}[Pendulum Period --- Key Properties]
\[
  T = 2\pi\sqrt{\frac{L}{g}}
\]
\begin{itemize}[noitemsep]
  \item $T$ depends on $L$ and $g$ only
  \item $T$ is \textbf{independent of mass}
  \item $T$ is \textbf{independent of amplitude}
  (for small angles)
  \item Longer pendulum $\implies$ longer period
  \item Stronger gravity $\implies$ shorter period
  \item A pendulum clock runs \textbf{slow} at higher
  altitude (smaller $g$)
\end{itemize}
\end{keybox}

\begin{diagbox}[The Simple Pendulum]
\begin{center}
\begin{tikzpicture}[scale=1.1]
  % Pivot
  \fill[gray!40] (-0.5,4) rectangle (0.5,4.3);
  \draw[thick] (-0.5,4) -- (0.5,4);
  \fill (0,4) circle (3pt);

  % String (equilibrium)
  \draw[thick, gray, dashed] (0,4) -- (0,1.5);

  % String (displaced)
  \draw[thick, codexblue] (0,4) -- (1.2,1.8);
  \fill[codexblue!50] (1.2,1.8) circle (8pt)
    node[right, black]{$m$};

  % Length label
  \draw[<->, gray] (-0.3,4) -- (-0.3,1.8);
  \node[left, gray] at (-0.3,2.9) {$L$};

  % Angle
  \draw[codexgold, thick] (0,3.2) arc(270:298:0.8);
  \node[codexgold] at (0.2,3.1) {$\theta$};

  % Forces on bob
  \draw[->, very thick, codexred]
    (1.2,1.8) -- (1.2,0.6) node[below]{$mg$};
  \draw[->, very thick, codexgreen]
    (1.2,1.8) -- (1.0,3.0) node[above left]{$T$};
  \draw[->, thick, codexblue, dashed]
    (1.2,1.8) -- (-0.2,1.8)
    node[left, font=\small, align=left]{$mg\sin\theta$\\(restoring)};

  % Period formula
  \node[align=center, codexblue, font=\small]
    at (-2.0, 0.5)
    {$T = 2\pi\sqrt{\dfrac{L}{g}}$};
\end{tikzpicture}
\end{center}
\end{diagbox}

\begin{examplebox}[7.2]
\stepcounter{workedexample}
A simple pendulum makes 20 complete oscillations in
$34\ \text{s}$.\\
(a) Find the period and frequency.\\
(b) Find the length of the pendulum.\\
(c) The pendulum is taken to a planet where
$g = 4.0\ \text{m\,s}^{-2}$.
What is the new period?
($g = 10\ \text{m\,s}^{-2}$ on Earth)
\end{examplebox}

\begin{solutionbox}
\textbf{(a)} Period and frequency:
\[
  T = \frac{34}{20} = 1.70\ \text{s}
  \qquad
  f = \frac{1}{T} = \frac{1}{1.70} = 0.588\ \text{Hz}
\]

\textbf{(b)} Length:
\[
  T = 2\pi\sqrt{\frac{L}{g}}
  \implies T^2 = 4\pi^2\frac{L}{g}
  \implies L = \frac{gT^2}{4\pi^2}
\]
\[
  L = \frac{10 \times 1.70^2}{4\pi^2}
  = \frac{10 \times 2.89}{39.48}
  = \frac{28.9}{39.48} = 0.732\ \text{m}
\]

\textbf{(c)} On the new planet:
\[
  T' = 2\pi\sqrt{\frac{L}{g'}}
  = 2\pi\sqrt{\frac{0.732}{4.0}}
  = 2\pi\sqrt{0.183}
  = 2\pi \times 0.428 = 2.69\ \text{s}
\]

The pendulum swings more slowly because gravity
is weaker.
\end{solutionbox}

% ============================================================
\section{The Mass-Spring System}
% ============================================================

\begin{processbox}[Period of a Mass-Spring System]
A mass $m$ on a spring of spring constant $k$ is
displaced by $x$ from equilibrium. By Hooke's Law,
the restoring force is:
\[
  F = -kx
\]
From Newton's second law $F = ma$:
\[
  a = -\frac{k}{m}x
\]
Comparing with $a = -\omega^2 x$: $\omega^2 = k/m$
\[
  \boxed{T = 2\pi\sqrt{\frac{m}{k}}}
  \qquad
  f = \frac{1}{2\pi}\sqrt{\frac{k}{m}}
\]
\end{processbox}

\begin{keybox}[Mass-Spring System --- Key Properties]
\[
  T = 2\pi\sqrt{\frac{m}{k}}
\]
\begin{itemize}[noitemsep]
  \item Heavier mass $\implies$ longer period (more
  inertia to oscillate)
  \item Stiffer spring (larger $k$) $\implies$ shorter
  period (stronger restoring force)
  \item \textbf{Period is independent of gravity} ---
  a spring clock works identically in space, on the
  Moon, and on Earth
  \item Period is independent of amplitude (for SHM
  conditions)
\end{itemize}
\end{keybox}

\begin{examplebox}[7.3]
\stepcounter{workedexample}
A $0.40\ \text{kg}$ mass hangs from a spring of
spring constant $100\ \text{N\,m}^{-1}$ and oscillates
vertically.\\
(a) Find the period and frequency.\\
(b) Find the amplitude if the maximum speed is
$0.60\ \text{m\,s}^{-1}$.\\
(c) Find the maximum acceleration.
\end{examplebox}

\begin{solutionbox}
\textbf{(a)} Period:
\[
  T = 2\pi\sqrt{\frac{m}{k}}
  = 2\pi\sqrt{\frac{0.40}{100}}
  = 2\pi\sqrt{0.004}
  = 2\pi \times 0.0632
  = 0.397\ \text{s}
\]
\[
  f = \frac{1}{T} = \frac{1}{0.397} = 2.52\ \text{Hz}
\]

\textbf{(b)} Angular frequency:
$\omega = 2\pi f = 2\pi \times 2.52 = 15.8\ \text{rad\,s}^{-1}$

From $v_{max} = A\omega$:
\[
  A = \frac{v_{max}}{\omega} = \frac{0.60}{15.8}
  = 0.038\ \text{m} = 3.8\ \text{cm}
\]

\textbf{(c)} Maximum acceleration:
\[
  a_{max} = A\omega^2 = 0.038 \times 15.8^2
  = 0.038 \times 249.6
  = 9.5\ \text{m\,s}^{-2}
\]
\end{solutionbox}

% ============================================================
\section{Energy in Simple Harmonic Motion}
% ============================================================

\begin{lawbox}[Energy in SHM]
In SHM, energy continuously converts between kinetic
and potential forms while total mechanical energy
remains constant.

\textbf{Kinetic energy:}
\[
  KE = \frac{1}{2}mv^2
  = \frac{1}{2}m\omega^2(A^2 - x^2)
\]

\textbf{Potential energy:}
\[
  PE = \frac{1}{2}m\omega^2 x^2
\]

\textbf{Total energy:}
\[
  E_{total} = KE + PE = \frac{1}{2}m\omega^2 A^2
  = \text{constant}
\]

At equilibrium ($x=0$): $KE = E_{total}$, $PE = 0$\\
At amplitude ($x=A$): $KE = 0$, $PE = E_{total}$
\end{lawbox}

\begin{diagbox}[Energy vs Displacement in SHM]
\begin{center}
\begin{tikzpicture}
\begin{axis}[
  width=10cm, height=6cm,
  xlabel={Displacement $x$},
  ylabel={Energy},
  xmin=-1.3, xmax=1.3,
  ymin=0, ymax=1.3,
  xtick={-1,0,1},
  xticklabels={$-A$,$0$,$A$},
  ytick={0,0.5,1},
  yticklabels={$0$,$\frac{1}{2}E$,$E_{total}$},
  grid=major, grid style={gray!20},
  legend pos=north east
]
  % Total energy (horizontal line)
  \addplot[black, thick, domain=-1.3:1.3]{1.0};
  \addlegendentry{$E_{total}$}

  % PE = (1/2)m*omega^2*x^2 (parabola)
  \addplot[codexred, thick, domain=-1.1:1.1,
    samples=50]{x^2};
  \addlegendentry{PE}

  % KE = E - PE
  \addplot[codexblue, thick, domain=-1:1,
    samples=50]{1 - x^2};
  \addlegendentry{KE}
\end{axis}
\end{tikzpicture}
\end{center}
\small At $x = A/\sqrt{2}$: KE = PE = $E_{total}/2$
\end{diagbox}

\begin{examplebox}[7.4]
\stepcounter{workedexample}
A particle of mass $0.20\ \text{kg}$ oscillates with
SHM of amplitude $0.10\ \text{m}$ and period
$1.5\ \text{s}$.\\
(a) Calculate the total energy.\\
(b) Find the displacement at which KE = PE.\\
(c) Find the KE and PE when $x = 0.06\ \text{m}$.
\end{examplebox}

\begin{solutionbox}
$\omega = 2\pi/T = 2\pi/1.5 = 4.19\ \text{rad\,s}^{-1}$

\textbf{(a)} Total energy:
\[
  E = \frac{1}{2}m\omega^2 A^2
  = \frac{1}{2}(0.20)(4.19)^2(0.10)^2
  = 0.1 \times 17.56 \times 0.01
  = 0.0176\ \text{J}
\]

\textbf{(b)} KE = PE means each $= E/2$:
\[
  PE = \frac{1}{2}m\omega^2 x^2 = \frac{E}{2}
\]
\[
  x^2 = \frac{A^2}{2} \implies x = \frac{A}{\sqrt{2}}
  = \frac{0.10}{\sqrt{2}} = 0.0707\ \text{m}
\]

\textbf{(c)} At $x = 0.06\ \text{m}$:
\[
  PE = \frac{1}{2}m\omega^2 x^2
  = \frac{1}{2}(0.20)(17.56)(0.06)^2
  = 0.1 \times 17.56 \times 0.0036
  = 0.00632\ \text{J}
\]
\[
  KE = E - PE = 0.0176 - 0.00632 = 0.0113\ \text{J}
\]
\end{solutionbox}

% ============================================================
\section{Damped and Forced Oscillations}
% ============================================================

\begin{definitionbox}[Damping]
\textbf{Damping} is the reduction in amplitude of
oscillation over time due to energy dissipation
(usually through friction or viscous drag).

\begin{itemize}[noitemsep]
  \item \textbf{Light damping:} Amplitude decreases
  gradually. The system oscillates many times before
  stopping. Example: a guitar string.
  \item \textbf{Heavy damping:} Amplitude decreases
  rapidly. The system returns slowly to equilibrium
  without oscillating. Example: a car shock absorber.
  \item \textbf{Critical damping:} The system returns
  to equilibrium as quickly as possible without
  oscillating. Example: door dampers, galvanometers.
  \item \textbf{Overdamping:} Returns to equilibrium
  slowly, without oscillating (but slower than
  critical damping).
\end{itemize}
\end{definitionbox}

\begin{diagbox}[Damping Comparison]
\begin{center}
\begin{tikzpicture}
\begin{axis}[
  width=12cm, height=6cm,
  xlabel={Time $t$},
  ylabel={Displacement $x$},
  xmin=0, xmax=12, ymin=-1.3, ymax=1.3,
  grid=major, grid style={gray!20},
  legend pos=north east,
  xtick=\empty, ytick=\empty
]
  % Light damping (oscillates, decays slowly)
  \addplot[codexblue, thick, domain=0:12, samples=200]
    {exp(-0.2*x)*sin(deg(2*x))};
  \addlegendentry{Light damping}

  % Heavy damping (slow return, no oscillation)
  \addplot[codexred, thick, domain=0:12, samples=100]
    {exp(-0.9*x)};
  \addlegendentry{Heavy damping}

  % Critical damping (fastest return without oscillation)
  \addplot[codexgreen, thick, domain=0:12, samples=100]
    {(1+3*x)*exp(-3*x)};
  \addlegendentry{Critical damping}

  % Equilibrium line
  \addplot[black, dashed, domain=0:12]{0};
\end{axis}
\end{tikzpicture}
\end{center}
\end{diagbox}

\begin{definitionbox}[Forced Oscillations and Resonance]
\textbf{Forced oscillations} occur when a periodic
external force drives a system to oscillate.
The system oscillates at the \textbf{driving frequency},
not its natural frequency.

\textbf{Resonance} occurs when the driving frequency
equals the \textbf{natural frequency} of the system.
At resonance:
\begin{itemize}[noitemsep]
  \item Amplitude of oscillation is maximum
  \item Energy transfer from driver to system
  is most efficient
  \item Can be destructive (Tacoma Narrows Bridge,
  1940 --- destroyed by wind-induced resonance)
  or useful (MRI machines, radio tuning, microwave
  ovens)
\end{itemize}
\end{definitionbox}

\begin{realworldbox}[Resonance in Bridges and Buildings]
On April 4, 2014, a suspension footbridge in Houphouet-Boigny,
Abidjan, Côte d'Ivoire, collapsed as crowds returning
from Easter mass marched in step across it. The rhythmic
footfall created a driving frequency close to the
bridge's natural frequency, producing resonance.
Engineers now design bridges to have natural
frequencies well away from typical human marching
frequencies, and instruct crowds to break step on
suspension bridges.
\end{realworldbox}

% ============================================================
\section{Chapter Summary}
% ============================================================

\begin{summarybox}
\textbf{SHM condition:} $a = -\omega^2 x$
(acceleration proportional to and opposite to
displacement)

\textbf{Equations:}
$x = A\sin(\omega t)$;
$v = A\omega\cos(\omega t) = \omega\sqrt{A^2-x^2}$;
$a = -\omega^2 x$

\textbf{Maxima:}
$v_{max} = A\omega$ (at $x=0$);
$a_{max} = A\omega^2$ (at $x=\pm A$)

\textbf{Simple pendulum:}
$T = 2\pi\sqrt{L/g}$
(independent of mass and amplitude)

\textbf{Mass-spring:}
$T = 2\pi\sqrt{m/k}$
(independent of amplitude and gravity)

\textbf{Energy:}
$E_{total} = \frac{1}{2}m\omega^2 A^2$;
$KE = \frac{1}{2}m\omega^2(A^2-x^2)$;
$PE = \frac{1}{2}m\omega^2 x^2$

\textbf{Damping:} reduces amplitude; critical damping
gives fastest return without oscillation.

\textbf{Resonance:} maximum amplitude when driving
frequency $=$ natural frequency.
\end{summarybox}

% ============================================================
\newpage
\begin{exercisebox}[7]

\subsection*{Section A: Multiple Choice}

\textbf{A1.} Which condition defines SHM?

\mcoption{A}{Constant velocity}
\mcoption{B}{Acceleration proportional to and in
the same direction as displacement}
\mcoption{C}{Acceleration proportional to and
opposite to displacement}
\mcoption{D}{Constant kinetic energy}
\ansref{Ch7-A1}

\vspace{0.4cm}

\textbf{A2.} A simple pendulum has period $T$ on Earth.
On a planet where $g$ is four times larger, the period is:

\mcoption{A}{$4T$}
\mcoption{B}{$2T$}
\mcoption{C}{$T/2$}
\mcoption{D}{$T/4$}
\ansref{Ch7-A2}

\vspace{0.4cm}

\textbf{A3.} In SHM, the velocity is maximum when:

\mcoption{A}{Displacement is maximum}
\mcoption{B}{Displacement is zero}
\mcoption{C}{Acceleration is maximum}
\mcoption{D}{Potential energy is maximum}
\ansref{Ch7-A3}

\vspace{0.4cm}

\textbf{A4.} A mass-spring system has period $T$.
If the mass is doubled and the spring constant is
halved, the new period is:

\mcoption{A}{$T$}
\mcoption{B}{$2T$}
\mcoption{C}{$T/2$}
\mcoption{D}{$T\sqrt{2}$}
\ansref{Ch7-A4}

\vspace{0.4cm}

\textbf{A5.} A particle oscillates with SHM of
amplitude $0.04\ \text{m}$ and angular frequency
$5\ \text{rad\,s}^{-1}$. Its maximum acceleration is:

\mcoption{A}{$0.20\ \text{m\,s}^{-2}$}
\mcoption{B}{$0.80\ \text{m\,s}^{-2}$}
\mcoption{C}{$1.00\ \text{m\,s}^{-2}$}
\mcoption{D}{$1.25\ \text{m\,s}^{-2}$}
\ansref{Ch7-A5}

\vspace{0.4cm}

\textbf{A6.} At which point in SHM are KE and PE equal?

\mcoption{A}{At the equilibrium position}
\mcoption{B}{At the amplitude}
\mcoption{C}{At $x = A/\sqrt{2}$}
\mcoption{D}{At $x = A/2$}
\ansref{Ch7-A6}

\vspace{0.4cm}

\textbf{A7.} Critical damping is best described as:

\mcoption{A}{The system oscillates with decreasing
amplitude}
\mcoption{B}{The system returns to equilibrium
as quickly as possible without oscillating}
\mcoption{C}{The system never returns to equilibrium}
\mcoption{D}{The amplitude remains constant}
\ansref{Ch7-A7}

\vspace{0.4cm}

\textbf{A8.} Resonance occurs when:

\mcoption{A}{The driving frequency is zero}
\mcoption{B}{The driving frequency equals the
natural frequency}
\mcoption{C}{Damping is maximum}
\mcoption{D}{The amplitude is minimum}
\ansref{Ch7-A8}

\vspace{0.4cm}

\textbf{A9.} A pendulum of length $1.0\ \text{m}$
has period (take $\pi^2 = 10$,
$g = 10\ \text{m\,s}^{-2}$):

\mcoption{A}{$1.0\ \text{s}$}
\mcoption{B}{$2.0\ \text{s}$}
\mcoption{C}{$\pi\ \text{s}$}
\mcoption{D}{$0.5\ \text{s}$}
\ansref{Ch7-A9}

\vspace{0.4cm}

\textbf{A10.} In SHM, the phase difference between
displacement and acceleration is:

\mcoption{A}{$0°$}
\mcoption{B}{$90°$}
\mcoption{C}{$180°$}
\mcoption{D}{$270°$}
\ansref{Ch7-A10}

\vspace{0.6cm}

\subsection*{Section B: Theory and Calculations}

\textbf{B1.} A particle oscillates with SHM given by
$x = 0.12\sin(8\pi t)$ metres, where $t$ is in seconds.\\[4pt]
(a) Write down the amplitude, angular frequency,
period and frequency.\\
(b) Find the maximum speed and maximum acceleration.\\
(c) Find the displacement, speed and acceleration
at $t = 0.10\ \text{s}$.\\
(d) At what displacement is the speed equal to half
the maximum speed?
\hfill\ansref{Ch7-B1}

\vspace{0.4cm}

\textbf{B2.} A student uses a simple pendulum to
determine $g$. She measures the following periods
for different lengths:

\begin{center}
\begin{tabular}{cc}
\toprule
$L$ (m) & $T$ (s)\\
\midrule
0.25 & 1.00\\
0.50 & 1.41\\
0.75 & 1.73\\
1.00 & 2.00\\
1.25 & 2.24\\
\bottomrule
\end{tabular}
\end{center}

(a) Explain why a graph of $T^2$ against $L$ gives
a straight line, and state the gradient in terms of $g$.\\
(b) Plot $T^2$ against $L$ (or use the data to find
the gradient algebraically) and determine $g$.\\
(c) The accepted value of $g$ is $9.81\ \text{m\,s}^{-2}$.
Calculate the percentage error in your answer.\\
(d) State two sources of experimental error in this
method.
\hfill\ansref{Ch7-B2}

\vspace{0.4cm}

\textbf{B3.} A $0.50\ \text{kg}$ mass hangs from
a spring. When pulled down $3.0\ \text{cm}$ from
equilibrium and released, it oscillates with period
$0.80\ \text{s}$.\\[4pt]
(a) Find the spring constant.\\
(b) Find the total energy of the oscillation.\\
(c) Find the maximum speed.\\
(d) Find the speed and KE when the displacement is
$2.0\ \text{cm}$.\\
(e) A second identical spring is attached in parallel
(i.e.\ both support the same mass). What is the new
period?
\hfill\ansref{Ch7-B3}

\vspace{0.4cm}

\textbf{B4.}
(a) Explain with a diagram what is meant by
(i) light damping, (ii) critical damping, and
(iii) heavy damping, giving one practical application
of each.\\[4pt]
(b) A car suspension spring has spring constant
$k = 25\,000\ \text{N\,m}^{-1}$. The car body it
supports has mass $m = 400\ \text{kg}$.\\
(i) Find the natural frequency of vertical oscillation.\\
(ii) Explain why the shock absorbers in a car should
produce critical (or near-critical) damping.\\
(iii) If the shock absorbers fail (no damping),
what happens when the car hits a pothole?
\hfill\ansref{Ch7-B4}

\vspace{0.4cm}

\textbf{B5.} A particle of mass $0.10\ \text{kg}$
undergoes SHM with amplitude $A = 0.15\ \text{m}$
and frequency $f = 3.0\ \text{Hz}$.\\[4pt]
(a) Calculate the total mechanical energy.\\
(b) Calculate the KE and PE at $x = 0.09\ \text{m}$.\\
(c) At what displacement is the KE equal to twice
the PE?\\
(d) A resistive force of $0.05\ \text{N}$ acts on
the particle throughout one complete oscillation.
Calculate the work done by this resistive force and
explain how this affects the subsequent motion.
\hfill\ansref{Ch7-B5}

\end{exercisebox}